\chapter{Zadatak 4c --- sinteza regulatora pomoću LMI-ja}
\label{ch:z4-lmi}

\section{Što se ovdje traži}

Naći matricu $K$ takvu da zatvoreni krug bude eksponencijalno stabilan
\emph{unatoč} prespajanju --- dakle riješiti upravo onaj problem koji je
podzadatak b) razotkrio.

\begin{zadatak}{Zadatak 4, podzadatak c)}
c) U sustavu (3), neka je matrica $B = \begin{bsmallmatrix} 1 \\ 0 \end{bsmallmatrix}$.
Dizajnirajte statički regulator $u(t) = Kx(t)$ takav da je zatvoreni krug
$\dot{x}(t) = (A(t) + BK)x(t)$ eksponencijalno stabilan, gdje je $A(t)$ zadana s
(4). Simulirajte ovaj vremenski promjenjiv zatvorni sustav s početnim uvjetom
$x(0) = \begin{bsmallmatrix} 1 \\ 0 \end{bsmallmatrix}$, prikažite $x(t)$ i
prikažite trajektoriju sustava u prostoru stanja.
\end{zadatak}

\begin{predznanje}
\begin{itemize}[nosep,leftmargin=*]
  \item Zašto Hurwitzovost modova nije dovoljna --- Poglavlje~\ref{ch:z4-simulacija}.
  \item Pozitivna definitnost i oznake $\pd$, $\nd$ --- \S\ref{sec:relacije}.
  \item Konveksna ljuska --- \S\ref{sec:z2-ljuska}.
  \item Inverz matrice --- \S\ref{sec:inverz}.
\end{itemize}
\end{predznanje}

\section{Ljapunovljev uvjet za jedan LTI sustav}
\label{sec:z4-ljapunov}

\begin{izvor}{Predavanje 7 --- Dinamički sustavi, Ljapunov, disipativnost, slajd 5/37 (PDF str.~6)}
\slajd[0.86]{p07-lti-stabilnost}

\textbf{Transkripcija.} \emph{Stabilnost vremenski invarijantnih sustava.}
Linearno vremenski invarijantan sustav $\dot{x}(t) = Ax(t)$ je eksponencijalno
stabilan \textbf{ako i samo ako} postoji $P$ tako da vrijedi
\[
P \pd 0 \quad \text{i} \quad A\trans P + PA \nd 0 .
\]
Ekvivalentan zapis gornje dvije nejednakosti:
$\begin{bsmallmatrix} -P & 0 \\ 0 & A\trans P + PA \end{bsmallmatrix} \nd 0$.
Ovo je linearna matrična nejednadžba (LMI) u $P$!

\textbf{Zašto je ovo baza cijelog podzadatka:} stabilnost se ovime prevodi iz
pitanja o svojstvenim vrijednostima u pitanje o \emph{dopustivosti} skupa
matričnih nejednakosti --- a to je konveksan problem koji rješavač može riješiti.
Uočite ,,ako i samo ako``: uvjet nije samo dovoljan nego i nužan.
\end{izvor}

\begin{izvor}{Predavanje 7 --- slajd 6/37 (PDF str.~10)}
\slajd[0.86]{p07-dokaz-trajektorije}

\textbf{Transkripcija (dokaz ,,putem trajektorija``).} Zbog stroge nejednakosti
$A\trans P + PA \nd 0$ postoji realan broj $\gamma > 0$ takav da vrijedi
$A\trans P + PA \nsd -\gamma P$. Za proizvoljnu trajektoriju:
\[
x\trans(A\trans P + PA)x \leq -\gamma\, x\trans P x
\;\Longrightarrow\;
\frac{\dd}{\dd t}\bigl(x\trans P x\bigr) \leq -\gamma\, x\trans P x ,
\]
odnosno, uz $V(t) = x(t)\trans P x(t)$, imamo $\dot{V} \leq -\gamma V$, pa slijedi
$V(t) \leq V(t_0)e^{-\gamma(t-t_0)}$. Uz poznate nejednakosti
$\lambda_{\min}(P)\norm{x}^2 \leq x\trans P x \leq \lambda_{\max}(P)\norm{x}^2$
slijedi
\[
\norm{x(t)} \leq \norm{x(t_0)}\sqrt{\frac{\lambda_{\max}(P)}{\lambda_{\min}(P)}}\;e^{-\frac{\gamma}{2}(t-t_0)} .
\]

\textbf{Odakle brojevi u našoj ocjeni.} Iz ovog izraza dolaze obje konstante koje
će se pojaviti u provjeri: stopa $\alpha = \gamma/2$ i predfaktor
$\sqrt{\lambda_{\max}(P)/\lambda_{\min}(P)} = \sqrt{\operatorname{cond}(P)}$.
Kod nas će biti $\alpha = 0{,}288374$ i predfaktor $1{,}3320$.
\end{izvor}

\begin{izvor}{Predavanje 7 --- slajd 7/37 (PDF str.~11)}
\slajd[0.86]{p07-nivo-skup}

\textbf{Transkripcija.} Iz $A\trans P + PA \nd 0$, množenjem s $x\trans$ slijeva i
$x$ zdesna:
\[
\dot{x}\trans P x + (Px)\trans \dot{x} < 0
\quad\Longleftrightarrow\quad
\dot{x}\trans \grad V(x) < 0, \qquad V(x) = x\trans P x, \ \grad V(x) = 2Px .
\]
Nivo-skup Ljapunovljeve funkcije $\{x \mid V(x) \leq c\}$ je \emph{invarijantan}
skup trajektorija sustava (jednom kad trajektorija uđe u taj skup, stalno ostaje
u njemu).

\textbf{Ovo je slika koju treba imati u glavi.} Nivo-skupovi funkcije
$V(x) = x\trans P x$ su elipse. Uvjet $\dot{V} < 0$ znači da trajektorija te
elipse siječe \emph{uvijek prema unutra} --- nikad ne može izaći. Upravo to će se
vidjeti na slici zatvorenog kruga.
\end{izvor}

\section{Kvadratna stabilnost: uvjet za \emph{skup} matrica}
\label{sec:z4-kvadratna}

\begin{izvor}{Predavanje 7 --- slajd 15/37 (PDF str.~25)}
\slajd[0.86]{p07-kvadratna-stabilnost}

\textbf{Transkripcija.} \emph{Robusna stabilnost, uvod.} Za zadani kompaktan
(zatvoren i ograničen) skup $\mathcal{A} \subset \R^{n\times n}$, promotrimo skup
dinamičkih sustava $\dot{x}(t) = Ax(t)$, $A \in \mathcal{A}$.

Kažemo da je skup sustava zadan s $\mathcal{A}$ \textbf{kvadratno stabilan} ako
postoji simetrična matrica $P$ takva da vrijedi
\[
P \pd 0, \qquad A\trans P + PA \nd 0 \quad \text{za sve } A \in \mathcal{A}.
\]
Funkcija $V(x) = x\trans P x$ je kvadratna Ljapunovljeva funkcija za sustave
$\mathcal{A}$ (stoga naziv ,,kvadratna stabilnost``).
\begin{itemize}[nosep,leftmargin=*]
  \item Implicira da su sve matrice $A \in \mathcal{A}$ stabilne (Hurwitzove).
  \item \emph{Implicira da za sve po dijelovima kontinuirane vremenski
        promjenjive $A(t) \in \mathcal{A}$, vremenski promjenjiv sustav
        $\dot{x}(t) = A(t)x(t)$ je eksponencijalno stabilan.}
\end{itemize}

\textbf{Druga natuknica je cijeli odgovor na Zadatak 4c.} Ključna riječ je
\textbf{zajednička}: jedna te ista matrica $P$ mora raditi za \emph{sve} modove.
Tada Ljapunovljeva funkcija pada bez obzira na to koji je mod trenutno aktivan,
pa prespajanje ne može naškoditi --- niti se mora znati \emph{kada} se prespaja.
Obrat ne vrijedi: postoje sklopni sustavi koji su stabilni, a nemaju zajedničku
kvadratnu Ljapunovljevu funkciju.
\end{izvor}

\begin{izvor}{Predavanje 7 --- slajd 16/37 (PDF str.~26)}
\slajd[0.86]{p07-konveksna-ljuska-lmi}

\textbf{Transkripcija.} Ako skup $\mathcal{A}$ ima beskonačno puno elemenata,
provjera kvadratne stabilnosti zahtijevala bi provjeru beskonačno puno LMI-jeva.
\emph{Ključno pitanje analize robusnosti: kako reducirati problem na standardan
LMI problem?}

Neka je $\mathcal{A}$ konveksna ljuska konačnog skupa matrica
$\{A_1, \dots, A_N\}$. \emph{Kvadratna stabilnost za konveksnu ljusku:} tada je
skup $\mathcal{A}$ kvadratno stabilan \textbf{ako i samo ako} postoji simetrična
$P$ takva da vrijedi
\[
P \pd 0, \qquad A_i\trans P + P A_i \nd 0 \quad \text{za sve } i = 1, \dots, N .
\]

\textbf{Ista ideja kao u Zadatku 2.} Ondje je beskonačan skup nejednakosti sveden
na vrhove pravokutnika; ovdje je beskonačan skup LMI-jeva sveden na generatore
konveksne ljuske. Kod nas je skup već konačan --- samo dvije matrice --- pa je
redukcija trivijalna: dovoljne su dvije nejednakosti.
\end{izvor}

\begin{formula}{Uvjet kvadratne stabilnosti za naš zatvoreni krug}
Uz $\mathcal{A} = \{A_{\mathrm{parno}} + BK,\ A_{\mathrm{neparno}} + BK\}$, tražimo
simetričnu $P$ takvu da
\begin{equation}
P \pd 0, \qquad \bar{A}\trans P + P\bar{A} \nd 0 \quad \text{za sve } \bar{A} \in \mathcal{A}.
\label{eq:z4-kvadratna}
\end{equation}
Tada je $V(x) = x\trans P x$ zajednička kvadratna Ljapunovljeva funkcija za oba
moda, a sustav je eksponencijalno stabilan za \emph{svaki} po dijelovima
kontinuiran izbor $\bar{A}(t) \in \mathcal{A}$ --- dakle i za zadani zakon
prespajanja.
\end{formula}

\section{Problem: uvjet nije LMI}
\label{sec:z4-nije-lmi}

\begin{izvor}{Predavanje 6 --- QCP i SDP, slajd 17/30 (PDF str.~17)}
\slajd[0.86]{p06-lmi-def}

\textbf{Transkripcija.} \emph{Linearna matrična nejednakost:}
\[
F(x) = F_0 + x_1F_1 + \dots + x_nF_n \nd 0 .
\]
\begin{itemize}[nosep,leftmargin=*]
  \item Vektor $x$ je vektor projektnih varijabli.
  \item Matrice $F_0, F_1, \dots, F_n$ su simetrične ($F_i = F_i\trans$), pa je i
        $F(x) = F(x)\trans$.
  \item $F(x)$ je (matrična) \emph{afina} funkcija od projektnih varijabli ---
        svaki element od $F(x)$ je afina funkcija od $x$.
  \item $\nd 0$ znači negativno definitno: $z\trans F(x) z < 0$ za svaki $z \neq 0$
        $\Leftrightarrow$ sve svojstvene vrijednosti od $F(x)$ su negativne
        $\Leftrightarrow$ $\lambda_{\max}(F(x)) < 0$.
\end{itemize}

\textbf{Riječ ,,afina`` je ovdje sve.} LMI je linearna matrična nejednakost
\emph{u projektnim varijablama}. Čim se u izrazu pojavi umnožak dviju
nepoznanica, to više nije LMI --- i upravo se to događa u sljedećem koraku.
\end{izvor}

\begin{korak}[Zašto uvjet \eqref{eq:z4-kvadratna} nije LMI]
Uvrštavanjem $\bar{A}_i = A_i + BK$:
\begin{equation}
(A_i + BK)\trans P + P(A_i + BK)
= A_i\trans P + \underbrace{K\trans B\trans P}_{\text{umnožak } K \text{ i } P}
+ PA_i + \underbrace{PBK}_{\text{umnožak } P \text{ i } K} \nd 0 .
\end{equation}
Nepoznanice su $P$ \emph{i} $K$, a pojavljuju se pomnožene jedna s drugom. Izraz
je \emph{bilinearan}, ne afin --- to je BMI (bilinearna matrična nejednakost), a
ne LMI. BMI je općenito nekonveksan problem i rješavači ga ne rješavaju
pouzdano.
\end{korak}

\section{Konveksifikacija: transformacija kongruencije}
\label{sec:z4-kongruencija}

\begin{izvor}{Predavanje 7 --- slajd 18/37 (PDF str.~32)}
\slajd[0.86]{p07-sinteza}

\textbf{Transkripcija.} \emph{Od analize do sinteze\dots{} ,,konveksifikacija``.}
Primjer --- statička povratna sprega varijabli stanja. Sustav
$\dot{x} = Ax + Bu$, regulator $u = Kx$; cilj: odrediti $K$ tako da je zatvoreni
krug stabilan. Zatvoreni krug: $\dot{x} = (A + BK)x$.

Uvjet stabilnosti zatvorenog kruga:
$P \pd 0$, $(A+BK)\trans P + P(A+BK) \nd 0$, odnosno
$P \pd 0$, $A\trans P + K\trans B\trans P + PA + PBK \nd 0$.
\emph{Nije LMI zbog umnoška varijabli $K$ i $P$!}

Nakon transformacije kongruencije s $P^{-1}$ (množenje slijeva s $(P^{-1})\trans$
i zdesna s $P^{-1}$) imamo (uočite $(P^{-1})\trans = P^{-1}$):
\[
P^{-1} \pd 0, \qquad P^{-1}A\trans + P^{-1}K\trans B\trans + AP^{-1} + BKP^{-1} \nd 0 .
\]
Uz supstituciju $\mathcal{Q} = P^{-1}$, $Y = KP^{-1}$ imamo:
\[
\mathcal{Q} \pd 0, \qquad \mathcal{Q}A\trans + Y\trans B\trans + A\mathcal{Q} + BY \nd 0
\qquad \textbf{LMI!}
\]

\textbf{Ovo je cijeli recept, doslovno.} Naš zadatak ga primjenjuje dvaput ---
jednom za svaki mod --- s istom $\mathcal{Q}$ i istim $Y$.
\end{izvor}

\begin{korak}[Zašto transformacija kongruencije čuva predznak]
Za invertibilnu $T$ i simetričnu $M$ vrijedi
\begin{equation}
M \nd 0 \iff T\trans M T \nd 0 .
\end{equation}
Dokaz je jedan redak: za proizvoljan $z \neq 0$ stavi $w = Tz$. Kako je $T$
invertibilna, $w \neq 0$, pa iz $M \nd 0$ slijedi
$z\trans (T\trans M T) z = w\trans M w < 0$. Obrat analogno, s $T^{-1}$.

Kod nas je $T = P^{-1}$, koja je invertibilna (jer je $P \pd 0$) i simetrična
(inverz simetrične matrice je simetričan), pa je $T\trans = T$. Množenje
nejednakosti s $P^{-1}$ s obje strane, dakle, ne mijenja njezin smisao --- za
razliku od množenja skalarom, gdje bi predznak bio problem.
\end{korak}

\begin{zamka}[Zamka: zašto se ne smije samo ,,podijeliti s $P$``]
Matrična nejednakost $M \nd 0$ nije jednadžba i ne smije se dijeliti. Jedina
dopuštena operacija koja čuva definitnost je kongruencija $M \mapsto T\trans M T$
s invertibilnom $T$. Ako se pokuša ,,skratiti`` $P$, dobiva se nešto što više
nije ekvivalentno polaznom uvjetu.
\end{zamka}

\begin{formula}{LMI za sintezu}
\begin{equation}
\mathcal{Q} \pd 0, \qquad
A_i \mathcal{Q} + \mathcal{Q}A_i\trans + BY + Y\trans B\trans \nd 0,
\qquad i \in \{\mathrm{parno},\ \mathrm{neparno}\},
\label{eq:z4-lmi}
\end{equation}
u nepoznanicama $\mathcal{Q} \in \Sym{2}$ i $Y \in \R^{1\times2}$. Iz rješenja se
rekonstruiraju
\begin{equation}
P = \mathcal{Q}^{-1}, \qquad K = Y P = Y\mathcal{Q}^{-1} .
\label{eq:z4-rekonstrukcija}
\end{equation}
\end{formula}

\begin{korak}[Broj nepoznanica]
$\mathcal{Q}$ je simetrična $2\times2$, dakle tri nezavisna broja
($q_{11}$, $q_{12}$, $q_{22}$); $Y$ je $1\times2$, dakle dva. Ukupno pet
skalarnih nepoznanica, i tri matrične nejednakosti. Za usporedbu, izravno
traženje $K$ (dva broja) i $P$ (tri broja) daje isti broj nepoznanica --- ali
nekonveksan problem. Supstitucija ne smanjuje veličinu problema, nego mu mijenja
\emph{klasu}.
\end{korak}

\begin{izvor}{Predavanje 6 --- QCP i SDP, slajd 20/30 (PDF str.~21)}
\slajd[0.86]{p06-lmi-konveksnost}

\textbf{Transkripcija.} Linearna matrična nejednadžba definira konveksna
ograničenja na $x$: skup $\mathcal{F} \coloneqq \{x \mid F(x) \nd 0\}$ je
konveksan skup. Isto vrijedi i za skupove s $\nsd$, $\pd$, $\psd$.
\emph{Dokaz:} za $x_1, x_2 \in \mathcal{F}$ i $\alpha \in [0,1]$ vrijedi
$F(\alpha x_1 + (1-\alpha)x_2) = \alpha F(x_1) + (1-\alpha)F(x_2) \nd 0$.
Funkcija $x \mapsto \lambda_{\max}(F(x))$ je konveksna funkcija.

\textbf{Zašto je to važno.} Ovo je razlog zbog kojeg se cijeli trud oko
konveksifikacije isplati: skup dopustivih $(\mathcal{Q}, Y)$ je konveksan, pa je
problem dopustivosti konveksan i rješavač ga rješava globalno --- nema lokalnih
rješenja u kojima bi mogao zapeti.
\end{izvor}

\section{Normalizacija strogih nejednakosti}
\label{sec:z4-normalizacija}

\begin{zamka}[Problem: rješavači ne znaju stroge nejednakosti]
Uvjeti \eqref{eq:z4-lmi} su \emph{strogi} ($\pd$, $\nd$), a numerički rješavači
rade s nestrogim ($\psd$, $\nsd$). Ako se strogo naivno zamijeni nestrogim,
dopustiv postaje i trivijalan par $\mathcal{Q} = 0$, $Y = 0$ --- iz kojeg se $K$
ne može rekonstruirati, jer $\mathcal{Q}$ nije invertibilna.
\end{zamka}

\begin{korak}[Rješenje: iskoristiti homogenost]
Sustav \eqref{eq:z4-lmi} je \textbf{homogen}: ako je $(\mathcal{Q}, Y)$ dopustiv,
dopustiv je i $(t\mathcal{Q},\ tY)$ za svaki $t > 0$, a pripadni regulator je
\emph{isti}:
\begin{equation}
K = (tY)(t\mathcal{Q})^{-1} = t Y \cdot \tfrac1t \mathcal{Q}^{-1} = Y\mathcal{Q}^{-1}.
\end{equation}
Rješenje se, dakle, smije proizvoljno skalirati. Iskoristi to i traži umjesto
strogih nejednakosti \emph{usklađeni par} normalizacija:
\begin{equation}
\mathcal{Q} \psd I, \qquad
A_i\mathcal{Q} + \mathcal{Q}A_i\trans + BY + Y\trans B\trans \nsd -I .
\label{eq:z4-normalizirano}
\end{equation}
Ovo su nestroge nejednakosti, pa ih rješavač prihvaća, a implicira ih strogost:
$\mathcal{Q} \psd I \Rightarrow \mathcal{Q} \pd 0$ i
$M_i \nsd -I \Rightarrow M_i \nd 0$. Trivijalno rješenje je isključeno.
\end{korak}

\section{Kod i rezultat}
\label{sec:z4-kod}

\begin{seminar}{Seminar v1.2, §5.4 --- ,,(c) Sinteza regulatora pomoću LMI-ja``}
Sustav je homogen: ako je $(\mathcal{Q}, Y)$ dopustiv, dopustiv je i
$(t\mathcal{Q}, tY)$ za $t \geq 1$, uz isti $K$. Stroge nejednakosti se stoga
numerički realiziraju usklađenim parom normalizacija $\mathcal{Q} \succeq I$ i
$A_i \mathcal{Q} + \mathcal{Q}A_i^{\mathsf{T}} + BY + Y^{\mathsf{T}}B^{\mathsf{T}} \preceq -I$,
a među dopustivim parovima bira se onaj s najmanjom normom $\|Y\|$. Problem se
rješava u YALMIP-u rješavačem SDPT3.

\begin{lstlisting}[style=matlabstyle]
Qv = sdpvar(2, 2, 'symmetric');
Yv = sdpvar(1, 2, 'full');
F  = [Qv >= eye(2), ...
      A_par*Qv + Qv*A_par' + B*Yv + Yv'*B' <= -eye(2), ...
      A_nep*Qv + Qv*A_nep' + B*Yv + Yv'*B' <= -eye(2)];
sol = optimize(F, norm(Yv, 2), sdpsettings('solver', 'sdpt3', 'verbose', 0));

Qs = value(Qv);  Ys = value(Yv);
P  = inv(Qs);    K  = Ys*P;
\end{lstlisting}

\begin{lstlisting}[style=outputstyle]
status : Successfully solved (SDPT3)
Q = [ 1.354304 -0.385757; -0.385757  1.420003]
Y = [-4.597618  1.172472]
P = [ 0.800314  0.217413;  0.217413  0.763286],  eig(P) = 0.563600, 1.000000
K = [-3.424629 -0.104650]
\end{lstlisting}
\end{seminar}

\begin{formula}{Rezultat sinteze}
\begin{equation}
P = \begin{bmatrix} 0{,}800314 & 0{,}217413 \\ 0{,}217413 & 0{,}763286 \end{bmatrix} \pd 0,
\qquad
K = \begin{bmatrix} -3{,}424629 & -0{,}104650 \end{bmatrix}.
\label{eq:z4-rjesenje}
\end{equation}
\end{formula}

\begin{zamka}[Nijansa: $\min\norm{Y}$ ne daje regulator najmanje norme]
Seminar kaže da se izborom najmanje $\norm{Y}$ ,,dobiva regulator najmanjeg
pojačanja``. To ne stoji: veza je $K = Y\mathcal{Q}^{-1}$, pa minimizacija
$\norm{Y}$ ne minimizira $\norm{K}$ --- između njih stoji $\mathcal{Q}^{-1}$,
koja je i sama varijabla. Rješavač smije smanjiti $\norm{Y}$ tako da istovremeno
,,napuše`` $\mathcal{Q}$, čime $\norm{K}$ može ostati velik.

Što ostaje istinito: dobiveni $K$ \emph{jest} valjan --- stabilizira zatvoreni
krug, i to je provjereno u Poglavlju~\ref{ch:z4-provjera}. Netočna je samo
dodatna tvrdnja o njegovoj optimalnosti.

\textbf{Kako to reći na obrani:} ,,Kriterij $\min\norm{Y}_2$ služio je samo kao
regularizacija --- da se među beskonačno mnogo dopustivih parova izabere jedan
određen i ponovljiv. Ne tvrdim da je dobiveni $K$ najmanje norme; za to bi
trebalo minimizirati $\norm{K}$ izravno, što uz $K = Y\mathcal{Q}^{-1}$ nije
konveksno u $(\mathcal{Q}, Y)$, ili nametnuti dodatno ograničenje oblika
$\mathcal{Q} \psd \eta I$ i $\begin{bsmallmatrix} \gamma I & Y \\ Y\trans & \mathcal{Q}\end{bsmallmatrix} \psd 0$.``
\end{zamka}

\begin{zamka}[Zamka: rješenje nije jedinstveno]
Sustav \eqref{eq:z4-lmi} je problem \emph{dopustivosti}, ne optimizacije. Drugi
rješavač, druga normalizacija ili drugi dodatni kriterij daju drugi $K$ koji
također stabilizira zatvoreni krug. To nije mana --- to je svojstvo problema. Na
obrani reci: ,,ovo je jedno od beskonačno mnogo dopustivih pojačanja; svi
rezultati i sve slike u radu dobiveni su s ovim.``
\end{zamka}

\begin{izgovor}
,,Traži se statički regulator koji stabilizira sustav unatoč prespajanju. Iz
podzadatka b) znam da Hurwitzovost pojedinih modova nije dovoljna, pa tražim
\emph{zajedničku} kvadratnu Ljapunovljevu funkciju --- jednu istu matricu pe koja
radi za oba moda. Prema Predavanju 7, ako takva postoji, sustav je eksponencijalno
stabilan za bilo koji zakon prespajanja, ne samo za zadani.

Uvjet je pe veće od nule i a s crtom transponirano pe plus pe a s crtom manje od
nule, za obje matrice zatvorenog kruga. Kad uvrstim a plus be ka, dobivam umnožak
nepoznanica ka i pe --- to je bilinearna, ne linearna matrična nejednakost, i
nije konveksno.

Rješenje je transformacija kongruencije s pe inverz: množim s obje strane, što
čuva predznak jer je pe inverz invertibilna i simetrična. Uz supstituciju ku
jednako pe inverz i ipsilon jednako ka pe inverz, sve postaje linearno u ku i
ipsilon --- dakle LMI.

Stroge nejednakosti sam numerički realizirao normalizacijama ku veće ili jednako
i, i izraz manji ili jednak minus i; to smijem jer je sustav homogen, pa
skaliranje ne mijenja ka.

Riješio sam u YALMIP-u sa SDPT3 i dobio ka jednako minus tri zarez četiri dva,
minus nula zarez jedan nula. Napomenuo bih da je ovo problem dopustivosti, pa
rješenje nije jedinstveno --- kriterij najmanje norme ipsilon služio je samo da
se izabere jedan određen par, a ne kao tvrdnja o optimalnosti pojačanja.``
\end{izgovor}
